\subsection{Primitive comparisons at arbitrary indices}
The next results distinguish Newton polygons without assuming that the absolute index is prime to $p$. Let $G,H$ be whole potentially supersingular height-two groups over a standing base $K$, with $p\ge5$, $\mu_p\subset K$ and $v=v_p(e_K)$. After completed unramified base change, use the common tame good field $L/K$ of degree $p^2-1$ and normalize $v_L(L^\times)=\Z$. A one-slope model has all nonzero $p$-point values equal to $e_K$. A two-slope model has an integer parameter $z$ with
\[
 0<z<pe_K,\qquad
 (p+1)e_K-z\quad(p-1\text{ points}),\qquad
 z/p\quad(p^2-p\text{ points}).
\]

\begin{theorem}[One-slope recovery]\label{thm:one-slope-all-indices}
If both good models have one first-level Newton slope and
\[
 n\ge v_p(e_K)+2,\qquad G[p^n]\simeq H[p^n],
\]
then $\theta_G=\theta_H$. There is no upper bound on $e_K$.
\end{theorem}

\begin{theorem}[An obstruction to different polygons]\label{thm:mixed-polygon-obstruction}
Let $n\ge1$. If $G[p^n]\simeq H[p^n]$, one model has one slope and the other has two slopes with parameter $z$, then
\begin{equation}\label{eq:mixed-polygon-obstruction}
 v_p(z)\ge\max\{1,2n-1-v_p(e_K)\},\qquad
 e_K>p^{\,2n-2-v_p(e_K)}.
\end{equation}
This assertion does not require unequal special characters.
\end{theorem}

\begin{corollary}[The remaining two-slope comparison]\label{cor:all-index-two-slope-reduction}
Let $p\ge5$, $n\ge2$, and let $M,M'$ belong to the full standing height-two filtration class over $K$. Suppose
\[
 e_K\le p^{\,2n-3-v_p(e_K)},\qquad M[p^n]\simeq M'[p^n].
\]
Their total dimensions and finite determinant factors agree. If their special characters differ, both groups are whole potentially supersingular groups and both good models have two first-level Newton slopes. In every other structural or polygon case the characters agree.
\end{corollary}

Define $\nu_{\mathrm{SS}}(p,e)$ to be the least nonnegative integer $r$ such that actual $p^{r+1}$-torsion recovers the special character uniformly among pairs of whole potentially supersingular height-two groups over the standing fields of index $e$.

\begin{theorem}[The full index and the supersingular index]\label{thm:full-supersingular-maximum}
For every $p\ge5$ and positive multiple $e$ of $p-1$, put $v=v_p(e)$. Then
\[
 \nuMin(p,e,2)=\max\left\{\nu_{\mathrm{SS}}(p,e),
       \left\lceil\frac{1+\log_p e+v}{2}\right\rceil\right\}.
\]
The index on the left concerns the full filtration class. The identity does not evaluate $\nu_{\mathrm{SS}}$ at every wild index.
\end{theorem}

\begin{theorem}[Reduction to a central comparison]\label{thm:all-index-central-reduction}
Under the hypotheses of Corollary~\ref{cor:all-index-two-slope-reduction}, suppose the special characters differ. For the two two-slope models one has
\[
 \max\{v_p(z_1),v_p(z_2)\}\ge v+2,\qquad v=v_p(e_K).
\]
Choose a model $G$ attaining this bound. There are $c$ successive quotients of $G$ by its canonical $p$-lines, paired with quotients of the other model by their transported lines, with
\[
 0\le c\le n-v-2,\qquad r=n-c\ge v+2,
\]
which give an unequal-character actual $p^r$-torsion comparison. The chosen terminal model has parameter $z_c=p^c z_G$ satisfying
\[
 e_K<z_c<pe_K,\qquad
 v_p\bigl((p+1)e_K-z_c\bigr)=v_p(e_K).
\]
The partner is whole potentially supersingular and may have either Newton polygon. The original numerical bound at level $n$ is not asserted at the smaller level $r$.
\end{theorem}

Section~\ref{sec:one-slope-all-indices} proves these assertions. The scaled torsion graph has $2n$ primitive Frobenius dimensions. A one-slope coordinate has a two-dimensional stable Frobenius span, while a two-slope coordinate has a one-dimensional span; this distinction controls overlaps of their nonzero weights. The primitive correspondence is classical \citep[\S\S2.1--2.2, pp.~3--4]{Koch2015}. The two-chain divisibility argument also proves Theorem~\ref{thm:all-index-upper}. At a failed canonical-line matching, Lemma~\ref{os-lem:canonical-matching} supplies an additional truncated valuation-sum inequality. For $v=0$, central $p^2$ rigidity excludes the terminal comparison; for $v>0$, the corresponding assertion with a $p$-divisible high numerator remains open.
